\section{Diferencias entre el muro de datos de los LLM y el desierto de datos de la robótica}

Aunque en ambos casos se dice «faltan datos», lo que les falta a los LLM y a los robots no es lo mismo. El problema de los LLM se parece más a un rendimiento marginal decreciente del corpus público de alta calidad: las páginas web, los libros, el código, los artículos científicos y los datos de preguntas y respuestas ya se han usado a gran escala, y seguir ampliándolos tropieza con cuellos de botella de repetición, contaminación, derechos de autor, calidad y evaluación. A la robótica, en cambio, le faltan desde el principio datos de interacción real: si no hay suficientes robots ejecutando tareas en entornos reales, no existe un flujo estable de datos desde los dispositivos.

\subsection{Dos tipos de escasez de datos}

\noindent\begin{longtable}{p{0.22\textwidth}p{0.35\textwidth}p{0.34\textwidth}}
\toprule
Objeto & Tipo de escasez de datos & Vías de solución habituales \\
\midrule
LLM & Rendimiento marginal decreciente del corpus de alta calidad; los benchmark se vuelven más difíciles & Mejor filtrado, mejores datos sintéticos, retroalimentación de orden superior, entornos de agent. \\
Robótica & Pocas trayectorias de acción reales, pocos escenarios de cola larga, recolección costosa & Simulación, teleoperación, despliegue real a pequeña escala, minería de datos de fallos. \\
Modelos del mundo & Datos insuficientes de predicción física y de dinámica espacial & Video, 3D, simulación, tareas de predicción multimodal. \\
VLA & Instrucciones, visión, acciones y retroalimentación difíciles de unificar & Trayectorias capturadas en el dispositivo, tareas de simulación, anotación en lenguaje, action head. \\
\bottomrule
\end{longtable}

\begin{importantbox}{El «muro de datos» no es un concepto unificado}
El muro de datos de los LLM es un problema de rendimiento marginal y de calidad; el desierto de datos de la robótica es un problema de interacción física y de escala de despliegue. Confundir estos dos conceptos conduce a estrategias de datos equivocadas.
\end{importantbox}

\subsection{Resumen de la sección}

Tanto a los LLM como a los robots les faltan datos, pero de manera distinta. Los primeros necesitan datos de mayor calidad y retroalimentación de orden superior; los segundos necesitan interacción física escalable y un ciclo cerrado de simulación.

